\chapter{Introducción al DEE}   % capítulos numerados normales
Este capítulo presenta el Drone Engineering Ecosystem (DEE), la plataforma académica desarrollada en la EETAC. El apartado 2.1 describe la arquitectura general del ecosistema y los tipos de módulos que lo componen. El apartado 2.2 presenta los drones utilizados en el DEE. El apartado 2.3 explica dónde se sitúa la contribución de este proyecto dentro del ecosistema.

\section{Arquitectura del DEE}

El DEE es una plataforma software abierta y modular que permite el control y la monitorización de uno o varios drones mediante diferentes tipos de dispositivos y aplicaciones. Su arquitectura se organiza en torno a cuatro tipos de módulos.

Los módulos a bordo se ejecutan en una Raspberry Pi montada sobre la plataforma del dron y se encargan de la interfaz directa con el autopilot, la cámara, los LEDs y otros periféricos. Para facilitar el despliegue, estos módulos se empaquetan en contenedores Docker.

Las aplicaciones frontend proporcionan la interfaz al operador e incluyen aplicaciones de escritorio (Python/TKinter, C\#), aplicaciones web (Flask, Vue.js) y aplicaciones móviles nativas (Android Studio, Flutter).

Los módulos backend se ocupan del almacenamiento y recuperación de datos.

La infraestructura de comunicación está formada por un broker interno y uno externo (público o alojado en la UPC), que conecta todos los módulos a través de internet. El modo de comunicación principal para las aplicaciones web y el despliegue remoto es el modo global, que enruta todos los mensajes a través del broker externo. El DEE soporta también un modo local, cuando no hay cobertura de internet, y un modo directo, empleado en simulación Software In The Loop (SITL), donde todos los módulos se ejecutan en un único dispositivo.


\section{Drones utilizados en el DEE}

El dron principal usado en el DEE consiste en un chasis HexSon con motores brushless y un controlador de vuelo Orange Cube PixHawk ejecutando firmware ArduPilot\cite{hexsoon}. La plataforma incorpora además una cámara. Este dron está destinado a operaciones en exteriores y es la plataforma de referencia para el control de todos los proyectos como EZDrone.

El DJI Tello es un dron pequeño, destinado a misiones en interiores. Dentro del DEE se utiliza principalmente en demostraciones a visitantes, dado que su tamaño reducido y su comportamiento estable lo hacen especialmente adecuado para entornos cerrados\cite{telloEETAC}. EZDrone incluye soporte para el control del Tello, permitiendo realizar estas demostraciones directamente desde el navegador.


\section{Contribución de este proyecto al DEE}
Este proyecto contribuye al bloque de aplicaciones Flutter del DEE con dos aportaciones complementarias. La primera es una colección de aplicaciones Flutter de demostración. Cada demostración está documentada de forma autónoma y está pensada para futuros desarrolladores. La segunda aportación (EZDrone), cae en webs de Flutter para el control de drones.

